\chapter{Methodology}

\section{Introduction}
This chapter describes the detection pipeline proposed in this thesis and the
systematic procedure used to build it. The design goal is a compact detector
that recovers very small human targets in aerial search-and-rescue imagery
without inflating the parameter or latency budget. Rather than adopting a single
fixed architecture, the model is developed through an \emph{additive ablation}:
starting from the YOLO11m baseline (reviewed in Chapter~II), one architectural
change is introduced at a time so that the contribution of each component can be
attributed independently. The chapter first presents the overall framework as a
structural anchor, then details each proposed modification, the state-space neck
variant that was explored, the ablation protocol, and a running example that
follows a single image through the pipeline. The baseline YOLO11 architecture and
the attention, multi-scale, and state-space concepts are covered in Chapter~II;
here the focus is strictly on the proposed modifications.

\section{Overall Framework}
The complete detection pipeline is shown in Figure~\ref{fig:framework}. An input
aerial image is letterbox-resized to $640 \times 640$ and passed through the
YOLO11m backbone, in which the terminal C2PSA attention block is replaced by a
Convolutional Block Attention Module (CBAM). The resulting multi-scale features
are aggregated by a PAN--FPN neck and decoded by a four-scale detection head that
adds a high-resolution P2 level (stride~4) to the standard P3--P5 levels. A final
non-maximum-suppression (NMS) stage produces the human bounding boxes. The two
proposed modifications, the CBAM attention block and the P2 detection scale,
are highlighted in the figure; the optional state-space neck variant
(Section~\ref{sec:mamba}) replaces selected neck blocks and is evaluated
separately.

\begin{figure}[tbp]
  \centering
  \figorplaceholder{figures/fig_overall_framework.png}{6cm}
  \caption{Overall framework of the proposed detector. An aerial image is
  preprocessed and passed through the YOLO11m backbone (with CBAM replacing the
  C2PSA block after the SPPF stage), a PAN--FPN neck, and a four-scale detection
  head; the added P2 level at stride~4 targets very small targets, and NMS yields
  the final human detections.}
  \label{fig:framework}
\end{figure}

\section{Detailed Methodology}

\subsection{Data Preprocessing}
All images are letterbox-resized to $640 \times 640$ so that the original aspect
ratio is preserved and no target is geometrically distorted, which matters when
the objects of interest are only a few pixels wide. Pixel values are normalised to
$[0,1]$, and only the default YOLO11 training-time augmentations are used
(mosaic, scaling, HSV jitter), with mosaic disabled for the final ten epochs
(\texttt{close\_mosaic}). Aggressive small-object augmentation (copy--paste,
mixup, vertical flip) was tested and discarded as a negative result, as reported
in Chapter~IV.

\subsection{Attention Enhancement: CBAM}
The baseline YOLO11m places two C2PSA self-attention blocks (1024 channels) after
the SPPF stage. In the proposed model these are replaced by a single CBAM module,
which applies channel attention followed by spatial attention in sequence. Given
an intermediate feature map $\mathbf{F} \in \mathbb{R}^{C \times H \times W}$, CBAM
computes a channel attention map $\mathbf{M}_c$ and a spatial attention map
$\mathbf{M}_s$,
\begin{align}
  \mathbf{M}_c(\mathbf{F}) &= \sigma\!\big(\mathrm{MLP}(\mathrm{AvgPool}(\mathbf{F}))
      + \mathrm{MLP}(\mathrm{MaxPool}(\mathbf{F}))\big), \\
  \mathbf{M}_s(\mathbf{F}') &= \sigma\!\big(f^{7\times7}\big([\mathrm{AvgPool}(\mathbf{F}');
      \mathrm{MaxPool}(\mathbf{F}')]\big)\big),
\end{align}
where $\sigma$ is the sigmoid function and $f^{7\times7}$ a convolution with a
$7\times7$ kernel. The features are refined multiplicatively,
$\mathbf{F}' = \mathbf{M}_c(\mathbf{F}) \otimes \mathbf{F}$ and
$\mathbf{F}'' = \mathbf{M}_s(\mathbf{F}') \otimes \mathbf{F}'$, where $\otimes$
denotes element-wise multiplication with broadcasting. Because CBAM is lighter
than the two C2PSA blocks it replaces, this substitution slightly \emph{reduces}
the parameter count relative to the baseline while adding both channel and spatial
selectivity.

\subsection{Multi-Scale Detection: the P2 Head}
The standard YOLO11 head detects at three scales: P3, P4 and P5 (strides 8, 16
and 32). For $640 \times 640$ input, the finest of these (P3) operates on an
$80 \times 80$ grid, so each cell already summarises an $8 \times 8$ pixel region;
targets smaller than this are difficult to localise. The proposed model adds a
fourth detection level, P2, at stride~4. This introduces a $160 \times 160$
feature map in which each cell covers only a $4 \times 4$ pixel region, allowing
the detector to resolve targets as small as approximately four pixels. The P2
level is formed by up-sampling and fusing the corresponding high-resolution
backbone features into the neck and attaching an additional detection branch, as
summarised in Figure~\ref{fig:framework}. As shown in Chapter~IV, this
high-resolution scale is the principal driver of the improvement in very-small-%
target recall.

\subsection{State-Space Neck Variant: C3k2Mamba}
\label{sec:mamba}
To test whether a selective state-space model (SSM) can improve feature
aggregation in the neck, an explored variant replaces the C3k2 bottleneck at six
neck layers (indices 13, 16, 19, 22, 25 and 28) with a C3k2Mamba block. Each such
block performs bidirectional local-window scanning, a forward and a reverse
selective scan over a local window (window sizes 6--8, state dimension
$d_{\text{state}} = 4$), and is inserted by post-initialisation surgical
replacement so that the surrounding CNN backbone and the P2 head are unchanged.
This variant is included for completeness and evaluated under the same protocol;
Chapter~IV reports that it does not improve accuracy while increasing both the
parameter count and the inference latency, and it is therefore excluded from the
recommended model.

\subsection{The Additive Ablation Design}
The four configurations evaluated in this thesis are summarised in
Table~\ref{tab:variants}. Each row introduces exactly one change with respect to
the preceding design, so that the effect of attention (CBAM), of the
high-resolution detection scale (P2), of their combination, and of the
state-space neck can each be isolated. All configurations are trained on the same
frozen C2A split with identical optimisation settings (AdamW, $\text{lr}_0 = 0.001$,
cosine schedule, up to 300 epochs with early stopping) and evaluated with the same
protocol and thresholds, so that every reported difference is attributable to the
architectural change alone. The experimental setup and these settings are detailed
in Chapter~IV.

\begin{table}[tbp]
  \centering
  \caption{The additive ablation. Each configuration differs from the previous one
  by a single architectural change, isolating the contribution of attention,
  multi-scale detection, and the state-space neck.}
  \label{tab:variants}
  \begin{tabular}{@{}l l l l@{}}
    \toprule
    Configuration & Backbone attention & Detection scales & Neck \\
    \midrule
    YOLO11m (baseline)   & C2PSA            & P3, P4, P5         & C3k2 \\
    + CBAM               & CBAM             & P3, P4, P5         & C3k2 \\
    + CBAM + P2          & CBAM             & P2, P3, P4, P5     & C3k2 \\
    + Mamba + CBAM + P2  & CBAM             & P2, P3, P4, P5     & C3k2Mamba (6 layers) \\
    \bottomrule
  \end{tabular}
\end{table}

\subsection{Running Example}
It helps to trace one image through the pipeline of Figure~\ref{fig:framework}.
Consider a collapsed-building scene of the kind shown in
Figure~\ref{fig:dataset_samples}(a), a $2000\times1500$ aerial photograph with more
than thirty people scattered across rubble, most of them under sixteen pixels tall.
The image is first letterbox-resized to $640\times640$, which preserves the aspect
ratio so the already small figures are not stretched. The backbone extracts features
at five scales, and the CBAM block after the SPPF stage re-weights them so that
human-shaped regions are emphasised over the surrounding debris. The neck fuses the
scales, and the head predicts boxes at all four levels; the added P2 level, at
stride four, is the one that fires on the smallest figures, which the stride-eight
P3 level would summarise into a single cell. Non-maximum suppression then removes
overlapping boxes, so that a cluster of nearby people is kept as separate detections
rather than merged. The output is one bounding box per recovered person, which is
what an operator reviews.

\section{Conclusion}
This chapter presented the proposed detector as an additive refinement of YOLO11m:
CBAM attention in place of C2PSA, a high-resolution P2 detection scale for very
small targets, and an explored state-space neck variant. The overall framework
anchors these modifications, and the additive ablation
design ensures that each component's contribution can be measured independently.
The next chapter reports the experimental setup and the quantitative and
qualitative results of this ablation.
